% Revised LaTeX replacement blocks for Chapter 5.
% Generated to align Chapter 5 with the frozen Chapter 4 methodology.
% These blocks are not compiled directly unless explicitly input from main.tex.

% -----------------------------------------------------------------------------
% BLOCK 1 -- Replace Experimental Evaluation Status table and introductory text
% -----------------------------------------------------------------------------

La valutazione sperimentale si basa sugli artefatti generati attraverso la
metodologia descritta nel Capitolo~\ref{chap:methodology}. La pipeline finale
comprende la costruzione del dataset congelato, la generazione degli split,
l'addestramento e la valutazione dei modelli proxy, la produzione delle
perturbazioni adversariali e delle trasformazioni anti-forensi, la costruzione
del forensic evaluation bundle e la normalizzazione degli output prodotti dagli
strumenti commerciali black-box. La valutazione commerciale è quindi trattata
come un livello operativo distinto rispetto ai proxy model, ma è inclusa nel
confronto finale mediante output osservabili e ricodificati rispetto alla ground
truth sperimentale \cite{sanna2024preliminary,sanna2025pixelperfect,magnet,cellebrite}.

Questa separazione è metodologicamente rilevante. Il livello dei modelli proxy
fornisce una baseline trasparente e riproducibile, poiché architetture,
checkpoint, fold di addestramento, predizioni e metriche sono disponibili per
ispezione e verifica. Il livello degli strumenti commerciali riflette invece una
condizione operativa differente: l'analista interagisce con sistemi professionali
o general-purpose la cui logica interna di classificazione non è pienamente
accessibile, e la valutazione deve quindi basarsi sugli output osservabili, sulle
categorie esportate, sui tag assegnati e sul matching rispetto al forensic
evaluation bundle \cite{casey2011digital,nist2006guide,iso27041,iso27042}.

\begin{table}[H]
\centering
\scriptsize
\renewcommand{\arraystretch}{1.15}
\setlength{\tabcolsep}{4pt}
\caption{Stato finale del workflow di valutazione sperimentale.}
\label{tab:results-evaluation-status}
\begin{tabularx}{\textwidth}{@{}p{0.28\textwidth} p{0.22\textwidth} X@{}}
\toprule
\textbf{Componente} & \textbf{Stato} & \textbf{Sintesi operativa} \\
\midrule
Dataset finale & \textbf{Congelato} & 1500 immagini bilanciate: 500 \texttt{weapon}, 500 \texttt{non\_weapon} e 500 \texttt{ood}. \\
Subset binario & \textbf{Congelato} & 1000 immagini: 500 \texttt{weapon} e 500 \texttt{non\_weapon}; usato per clean, adversarial e anti-forensic evaluation. \\
Split clean e \gls{ood} & \textbf{Completati} & Cinque fold binari clean e un insieme \gls{ood} separato, derivati dai manifest congelati. \\
Addestramento proxy & \textbf{Completato} & Checkpoint fold-aware prodotti per \texttt{efficientnet\_b0}, \texttt{resnet18} e \texttt{clip}. \\
Valutazione proxy & \textbf{Completata} & Valutazione su input clean, \gls{ood}, adversariali e anti-forensi. \\
Attacchi adversariali & \textbf{Generati} & \gls{fgsm}, \gls{superdeepfool}, \gls{sigmazero}, \gls{onepixel} e \gls{colorshift}. \\
Trasformazioni anti-forensi & \textbf{Generate} & \gls{jpegrecompression}, \gls{resampleresize}, \gls{gaussianblur}, \gls{histogrammodification} e \gls{contraststretching}. \\
Forensic evaluation bundle & \textbf{Generato} & 11\,500 file predisposti per la valutazione dei proxy model e degli strumenti black-box. \\
\gls{magnetaxiomtool} / \gls{magnetai} & \textbf{Consolidato} & Bundle processato ed export normalizzato; il tag \texttt{Possible weapons} è ricondotto alla predizione positiva \texttt{weapon}. \\
\gls{excirefoto2025} & \textbf{Consolidato} & Valutato come software standalone general-purpose di semantic retrieval, con tre configurazioni operative: \texttt{D20}, \texttt{D50}, \texttt{D80}. \\
\gls{cellebriteinseyets} & \textbf{Consolidato} & Export normalizzato da ambiente \texttt{Cellebrite Physical Analyzer 10.9.0.3029}; le classificazioni osservabili sono ricodificate nel task operativo \texttt{weapon}/\texttt{non\_weapon}. \\
Explainability & \textbf{Integrata} & Casi rappresentativi generati con \gls{ig}; analisi qualitativa e diagnostica, non metrica primaria. \\
\bottomrule
\end{tabularx}
\end{table}

% -----------------------------------------------------------------------------
% BLOCK 2 -- Replace Excire subsection title and opening paragraphs
% -----------------------------------------------------------------------------

\subsection{\gls{excirefoto2025} results and sensitivity analysis}
\label{subsec:results-excire}

La valutazione di \gls{excirefoto2025} richiede una cautela interpretativa
specifica. Il software non viene trattato come prodotto sviluppato esclusivamente
per finalità forensi, né come classificatore nativo per la distinzione binaria
\texttt{weapon}/\texttt{non\_weapon}. Nel protocollo \gls{fairlab}, esso viene
valutato come sistema black-box di semantic retrieval con funzionalità
\gls{ai}-based. La predizione operativa viene quindi derivata dalla presenza o
assenza del campione tra i risultati recuperati rispetto a prompt firearm-oriented
e a una specifica configurazione del parametro \texttt{Distance limit}. Questa
ricodifica consente una valutazione quantitativa rispetto alla ground truth
sperimentale, ma non deve essere interpretata come accesso alla logica
classificatoria interna del sistema.

Per valutare la sensibilità del comportamento operativo alla soglia di recupero,
sono state considerate tre configurazioni del parametro \texttt{Distance limit}:
\texttt{D20}, \texttt{D50} e \texttt{D80}. La configurazione \texttt{D20}
rappresenta un'impostazione restrittiva, orientata a ridurre falsi positivi e
recuperi semanticamente deboli; \texttt{D50} costituisce la configurazione
intermedia della sensitivity analysis; \texttt{D80} rappresenta invece
un'impostazione permissiva, orientata a massimizzare il recall ma esposta a un
aumento rilevante dei falsi positivi.

% -----------------------------------------------------------------------------
% BLOCK 3 -- New Cellebrite Inseyets subsection
% -----------------------------------------------------------------------------

\subsection{\gls{cellebriteinseyets} results}
\label{subsec:results-cellebrite}

La valutazione di \gls{cellebriteinseyets} completa il livello commerciale
black-box della pipeline. Il tool viene considerato come piattaforma commerciale
di analisi assistita dei contenuti multimediali, valutata attraverso gli output
osservabili esportati nell'ambiente \texttt{Cellebrite Physical Analyzer
10.9.0.3029}. In coerenza con la metodologia del Capitolo~\ref{chap:methodology},
\gls{ufed} non viene trattato come classificatore immagini, ma solo come eventuale
componente dell'ecosistema Cellebrite di acquisizione e gestione dei dati.

\begin{table}[H]
\centering
\caption{Prestazioni complessive di \gls{cellebriteinseyets} sul forensic evaluation bundle.}
\label{tab:results-cellebrite-global-metrics}
\small
\begin{tabular}{lr}
\toprule
\textbf{Metrica} & \textbf{Valore} \\
\midrule
Accuracy complessiva sul task binario & 0.958 \\
Balanced accuracy & 0.958 \\
Precision \texttt{weapon} & 0.958 \\
Recall \texttt{weapon} & 0.958 \\
F1 \texttt{weapon} & 0.958 \\
False negative rate & 0.042 \\
False positive rate & 0.042 \\
True positives & 5\,268 \\
True negatives & 5\,271 \\
False positives & 229 \\
False negatives & 232 \\
\gls{ood} weapon flag rate & 0.292 \\
\bottomrule
\end{tabular}
\end{table}

Nel confronto aggregato, \gls{cellebriteinseyets} mostra il profilo commerciale
più stabile tra i tool finali: l'accuracy complessiva sul task binario è pari a
0.958, con \gls{fnr} e \gls{fpr} entrambi prossimi a 0.042. Tale equilibrio è
operativamente rilevante perché indica una distribuzione relativamente bilanciata
tra rischio di mancata segnalazione della classe \texttt{weapon} e rischio di
rumore operativo generato da falsi positivi. Il valore non deve tuttavia essere
interpretato come misura della robustezza interna del prodotto, poiché la logica
di classificazione rimane black-box e l'analisi si basa esclusivamente sugli
output esportati e normalizzati.

\begin{table}[H]
\centering
\caption{Comportamento di \gls{cellebriteinseyets} per condizione sperimentale.}
\label{tab:results-cellebrite-condition-metrics}
\small
\begin{tabularx}{\textwidth}{@{}l r r r r X@{}}
\toprule
\textbf{Condizione} & \textbf{Accuracy} & \textbf{Recall weapon} & \textbf{\gls{fnr}} & \textbf{\gls{fpr}} & \textbf{Nota operativa} \\
\midrule
Clean & 0.964 & 0.968 & 0.032 & 0.040 & Baseline commerciale su 1\,000 immagini non perturbate. \\
Perturbed & 0.958 & 0.957 & 0.043 & 0.042 & Valori calcolati su 10\,000 immagini perturbate: 5\,000 adversariali e 5\,000 anti-forensi. \\
Adversarial & 0.955 & 0.950 & 0.050 & 0.040 & Degrado contenuto rispetto alla baseline clean; \gls{fgsm} rappresenta la condizione più critica. \\
Anti-forensic & 0.960 & 0.963 & 0.037 & 0.043 & Stabilità elevata rispetto alle trasformazioni model-agnostic considerate. \\
\gls{ood} & -- & -- & -- & -- & Weapon flag rate pari a 0.292: 146 immagini \gls{ood} su 500 sono state ricondotte alla classe operativa \texttt{weapon}. \\
\bottomrule
\end{tabularx}
\end{table}

Il comportamento sulle immagini \gls{ood} è più contenuto rispetto a quello
osservato per \gls{clip}, \gls{magnetai} ed \gls{excirefoto2025} nella
configurazione \texttt{D50}. In particolare, 146 immagini \gls{ood} su 500 sono
ricondotte alla classe operativa \texttt{weapon}. Questo valore non rappresenta
un errore binario in senso stretto, poiché i campioni \gls{ood} non appartengono
alla ground truth \texttt{weapon}/\texttt{non\_weapon}; esso descrive piuttosto
la tendenza del tool a segnalare come potenzialmente rilevanti contenuti
borderline, ambigui o semanticamente weapon-adjacent.

\begin{table}[H]
\centering
\caption{Prestazioni di \gls{cellebriteinseyets} nelle condizioni adversariali e anti-forensi.}
\label{tab:results-cellebrite-attack-metrics}
\scriptsize
\renewcommand{\arraystretch}{1.12}
\setlength{\tabcolsep}{4pt}
\begin{tabular}{@{}l l r r r r r r@{}}
\toprule
\textbf{Famiglia} & \textbf{Condizione} & \textbf{Accuracy} & \textbf{Recall} & \textbf{\gls{fnr}} & \textbf{\gls{fpr}} & \textbf{FN} & \textbf{FP} \\
\midrule
Adversarial & \gls{colorshift} & 0.964 & 0.968 & 0.032 & 0.040 & 16 & 20 \\
Adversarial & \gls{fgsm} & 0.918 & 0.894 & 0.106 & 0.058 & 53 & 29 \\
Adversarial & \gls{onepixel} & 0.962 & 0.968 & 0.032 & 0.044 & 16 & 22 \\
Adversarial & \gls{sigmazero} & 0.970 & 0.958 & 0.042 & 0.018 & 21 & 9 \\
Adversarial & \gls{superdeepfool} & 0.961 & 0.964 & 0.036 & 0.042 & 18 & 21 \\
\midrule
Anti-forensic & \gls{contraststretching} & 0.967 & 0.972 & 0.028 & 0.038 & 14 & 19 \\
Anti-forensic & \gls{gaussianblur} & 0.954 & 0.964 & 0.036 & 0.056 & 18 & 28 \\
Anti-forensic & \gls{histogrammodification} & 0.949 & 0.940 & 0.060 & 0.042 & 30 & 21 \\
Anti-forensic & \gls{jpegrecompression} & 0.966 & 0.972 & 0.028 & 0.040 & 14 & 20 \\
Anti-forensic & \gls{resampleresize} & 0.964 & 0.968 & 0.032 & 0.040 & 16 & 20 \\
\bottomrule
\end{tabular}
\end{table}

La condizione più critica per \gls{cellebriteinseyets} è \gls{fgsm}, con accuracy
pari a 0.918 e \gls{fnr} pari a 0.106. A differenza di quanto osservato per il
proxy target EfficientNet-B0 sotto \gls{sigmazero}, il tool non presenta una
compromissione estrema in corrispondenza degli attacchi sparse o decision-boundary
considerati. Tra le trasformazioni anti-forensi, \gls{histogrammodification}
produce il valore di \gls{fnr} più alto, pari a 0.060, mentre \gls{gaussianblur}
produce il valore di \gls{fpr} più elevato, pari a 0.056. Nel complesso, il
profilo osservato suggerisce una maggiore stabilità operativa rispetto alle
perturbazioni del bundle, pur con falsi negativi residui rilevanti per un
workflow di triage.

% -----------------------------------------------------------------------------
% BLOCK 4 -- Replace Comparative Robustness Discussion opening and summary table
% -----------------------------------------------------------------------------

\section{Comparative Robustness Discussion}
\label{sec:results-comparative}

I risultati ottenuti consentono di confrontare tre livelli della pipeline:
modelli proxy trasparenti, perturbazioni sperimentali controllate e valutazione
black-box commerciale. I modelli proxy permettono di osservare in modo
controllato il degrado prestazionale rispetto a clean, \gls{ood}, adversarial e
anti-forensic input. \gls{magnetai}, \gls{excirefoto2025} e
\gls{cellebriteinseyets} permettono invece di verificare come strumenti
commerciali black-box reagiscano allo stesso bundle sperimentale, pur senza
accesso alla logica interna dei rispettivi classificatori o motori di retrieval.

\begin{table}[H]
\centering
\caption{Sintesi comparativa tra modelli proxy e strumenti commerciali black-box.}
\label{tab:results-comparative-summary}
\small
\begin{tabularx}{\textwidth}{@{}p{2.8cm} r r r X@{}}
\toprule
\textbf{Sistema} & \textbf{Clean accuracy} & \textbf{Clean \gls{fnr}} & \textbf{OOD weapon rate} & \textbf{Osservazione principale} \\
\midrule
EfficientNet-B0 & 0.978 & 0.016 & 0.370 & Migliore baseline clean tra i proxy, ma vulnerabilità marcata ad alcune perturbazioni adversariali, in particolare \gls{sigmazero} e \gls{fgsm}. \\
ResNet18 & 0.964 & 0.030 & 0.438 & Buona baseline e maggiore stabilità rispetto a EfficientNet-B0 per alcune perturbazioni generate nello scenario white-box considerato. \\
\gls{clip} & 0.927 & 0.012 & 0.836 & Accuracy clean inferiore ai modelli supervisionati, ma forte tendenza a segnalare come \texttt{weapon} immagini fuori distribuzione. \\
\gls{magnetaxiomtool} / \gls{magnetai} & 0.950 & 0.070 & 0.360 & Comportamento relativamente conservativo: basso \gls{fpr}, ma rischio non trascurabile di falsi negativi sulla classe \texttt{weapon}. \\
\gls{excirefoto2025} \texttt{D50} & 0.944 & 0.042 & 0.340 & Configurazione intermedia di triage: recall elevato sulla classe \texttt{weapon}, con incremento controllato dei falsi positivi rispetto a \texttt{D80}. \\
\gls{cellebriteinseyets} & 0.964 & 0.032 & 0.292 & Profilo commerciale più stabile nel bundle finale, con buon equilibrio tra falsi negativi, falsi positivi e weapon flag rate sugli input \gls{ood}. \\
\bottomrule
\end{tabularx}
\end{table}

Il confronto non deve essere interpretato come ranking assoluto tra modelli proxy
e strumenti commerciali, perché i primi espongono predizioni controllate su classi
binarie note, mentre i secondi operano tramite categorie proprietarie, tag
esportati, classificazioni osservabili o recupero semantico. Il valore del
confronto è operativo: mostrare se le tendenze di fragilità osservate nei proxy si
riflettono anche in strumenti black-box utilizzati in workflow investigativi.

La presenza di \gls{cellebriteinseyets} modifica la lettura del livello
commerciale. Rispetto a \gls{magnetai} ed \gls{excirefoto2025} nella
configurazione \texttt{D50}, \gls{cellebriteinseyets} mostra una clean accuracy
pari a 0.964, un clean \gls{fnr} pari a 0.032 e il weapon flag rate \gls{ood} più
contenuto tra gli strumenti commerciali finali. Questo non consente di attribuire
al tool una robustezza generale o certificata, ma documenta un comportamento
operativo più stabile nel bundle sperimentale considerato.

La differenza tra i sistemi conferma che la robustezza non può essere dedotta
dalla sola accuracy clean, ma deve essere valutata rispetto a condizioni
specifiche di impiego, inclusi input anomali, borderline o semanticamente
distanti dalla distribuzione attesa. Nei workflow di \gls{df}, il problema non è
soltanto stabilire quale sistema sia più accurato in condizioni nominali, ma
comprendere quale profilo di errore esso introduca nel processo di triage:
mancata segnalazione di contenuti rilevanti, aumento del rumore operativo,
assorbimento di casi \gls{ood} o sensibilità a perturbazioni mirate.
